\documentclass[10pt]{article}
\usepackage[T1]{fontenc}
\usepackage[utf8]{inputenc}
\usepackage[french]{babel}
\usepackage{amsmath,amssymb,mathrsfs,mathpazo}
\usepackage{geometry}
\geometry{a4paper,margin=2cm}
\usepackage{graphicx,array,tabularx,tabularray,colortbl,multirow,multicol,booktabs,enumitem,pifont,ulem,fancybox,diagbox,cancel}
\usepackage{pstricks,pstricks-add,pst-plot,pst-tree,pst-node,pst-text,pst-func,pst-math,pst-bspline,pst-3dplot,pst-eucl,pst-solides3d}
\let\pscalculate\undefined
\usepackage{pst-bezier}
\usepackage{tikz,pgfplots,tkz-tab}
\usetikzlibrary{arrows,arrows.meta,patterns,calc,positioning,trees,decorations.markings,shapes.geometric,shapes.misc,shapes.arrows,intersections,decorations.pathmorphing,decorations.pathreplacing,backgrounds,shadows,babel}
\pgfplotsset{compat=1.18}
\usepgfplotslibrary{fillbetween}
\usepackage[np]{numprint}
\usepackage{listings,makecell,pas-tableur,twoopt}
\let\barre\undefined
\usepackage[pstricks]{tablvar}
\definecolor{darkgreen}{RGB}{0,100,0}
\providecommand{\red}{\color{red}}
\providecommand{\blue}{\color{blue}}
\providecommand{\green}{\color{green}}
\providecommand{\black}{\color{black}}
\providecommand{\white}{\color{white}}
\providecommand{\gray}{\color{gray}}
\providecommand{\cyan}{\color{cyan}}
\providecommand{\magenta}{\color{magenta}}
\providecommand{\yellow}{\color{yellow}}
\providecommand{\violet}{\color{violet}}
\providecommand{\orange}{\color{orange}}
\providecommand{\euro}{\texteuro}
\providecommand{\EUR}{\texteuro}
\providecommand{\ieme}{\textsuperscript{e}}
\providecommand{\degres}{\ensuremath{^\circ}}
\providecommand{\degree}{\ensuremath{^\circ}}
\providecommand{\vv}[1]{\overrightarrow{#1}}
\providecommand{\Coord}[1]{#1}
\providecommand{\decosix}{}
\providecommand{\starredbullet}{\ensuremath{\bullet}}
\pagestyle{empty}
\setlength{\parindent}{0pt}
\usepackage{adjustbox}\DeclareUnicodeCharacter{2236}{:}\DeclareUnicodeCharacter{2190}{\ensuremath{\leftarrow}}
\providecommand{\R}{\mathbb{R}}
\providecommand{\N}{\mathbb{N}}
\providecommand{\Z}{\mathbb{Z}}
\providecommand{\Q}{\mathbb{Q}}
\providecommand{\D}{\mathbb{D}}
\providecommand{\C}{\mathbb{C}}
\providecommand{\e}{{\mathrm{e}}}
\providecommand{\vect}[1]{\overrightarrow{#1}}
\providecommand{\vectt}[1]{\overrightarrow{\mathrm{#1}}}
\providecommand{\vvt}[1]{\overrightarrow{\mathrm{#1}}}
\providecommand{\np}[2][]{{#2~#1}}
\providecommand{\vv}[1]{{\overrightarrow{#1}}}
\renewcommand{\d}{{\mathrm{d}}}
\providecommand{\degres}{{^\circ}}
\providecommand{\degre}{{^\circ}}
\providecommand{\blue}{\color{blue}}
\providecommand{\red}{\color{red}}
\providecommand{\enskip}{\quad}
\providecommand{\llbracket}{\lbrack\!\lbrack}
\providecommand{\rrbracket}{\rbrack\!\rbrack}
\providecommand{\euro}{{\text{€}}}
\providecommand{\no}{{\text{n°}}}
\providecommand{\interval}[2]{{[#1;#2]}}
\providecommand{\intervFO}[2]{{[#1;#2[}}
\providecommand{\dotfill}{\ldots}

\providecommand{\ds}{\displaystyle}
\providecommand{\textding}[1]{#1}
\providecommand{\Oij}{\(\mathrm{O};\vec{\imath},\vec{\jmath}\)}
\providecommand{\Oijk}{\(\mathrm{O};\vec{\imath},\vec{\jmath},\vec{k}\)}
\providecommand{\Ouv}{\(\mathrm{O};\vec{u},\vec{v}\)}
\providecommand{\intervFF}[2]{\ensuremath{[#1;#2]}}
\renewcommand{\euro}{\ifmmode\text{\texteuro}\else\texteuro\fi}
\DeclareUnicodeCharacter{00B0}{\ifmmode^{\circ}\else\textdegree\fi}
\DeclareUnicodeCharacter{20AC}{\ifmmode\text{\texteuro}\else\texteuro\fi}
% Énoncés : Éducation nationale. Transcription et corrigés : APMEP, Denis Vergès et contributeurs. Découpage et classement : CoopMaths, Sébastien Lozano, Éric Elter et contributeurs. Adaptation de la présentation : Maths Explorer.
% CC BY-SA : https://coopmaths.fr/www/about/
% https://coopmaths.fr/alea/static/bac/2023/tex/bac_2023_09_sujet1_ameriquesud_4_cor.tex
\begin{document}
\noindent\textit{APMEP / CoopMaths. CC BY-SA. Présentation adaptée par Maths Explorer.}\par\medskip
\medskip

\textbf{Partie A}

\medskip


\[u_{n+1} =  2u_n\left(1 - u_n\right).\]

Cette relation de récurrence s'écrit $u_{n+1} =  f\left(u_n\right)$, où $f$ est la fonction définie sur $\R$ par :
\[f(x) = 2x(1 - x).\]

\smallskip
\begin{enumerate}
\item 
$f$ est une fonction polynomiale donc dérivable sur $\R$ et sur cet intervalle :

$f(x) = 2x - 2x^2$, d'où $f'(x) = 2 - 4x = 2(1 - 2x)$.

On a :

$\bullet~~$$f'(x) = 0 \iff 1 - 2x = 0 \iff x = \dfrac12$ ;

$\bullet~~$$f'(x) > 0 \iff 1 - 2x > 0 \iff x < \dfrac12$ ; la dérivée est positive sur $\left[0~;~\dfrac12\right]$, la fonction $f$ est donc croissante sur cet intervalle.

$\bullet~~$$f'(x) < 0 \iff 1 - 2x < 0 \iff x > \dfrac12$ ;
\item On admet que pour tout entier naturel $n,\:0 \leqslant u_n \leqslant \dfrac12$.


$\bullet~~$$u_1 = 2u_0\left(1 - u_0\right) = 2\times 0,3 \times 0,7= 0,42$.

$\bullet~~$Démonstration par récurrence :

\begin{itemize}
\item \emph{Initialisation} : on a $u_0 \leqslant u_1$ car $0,3 \leqslant 0,42$.
\item \emph{Hérédité} On suppose qu'il existe $n \in \N$  tel que $u_n \leqslant u_{n+1}$.

Comme on suppose que $u_n \leqslant u_{n+1} \leqslant \dfrac12$ on a par croissance de la fonction $f$ sur $\left[0~;~\dfrac12\right], \: \:f\left(u_n\right) \leqslant f\left(u_{n+1}\right) \leqslant f\left(\dfrac12\right)$, soit 

$u_{n+1} \leqslant u_{n+2} \leqslant 2 \times \dfrac12 \times \dfrac12$, ou encore $u_{n+1} \leqslant u_{n+2} \leqslant \dfrac12$ : l'encadrement est vrai au rang $n+1$.

La relation est vraie au rang 0 et si elle est vraie au rang $n \in \N$ elle l'est encore au rang $n + 1$ : d'après le principe de récurrence on a démontré que :

pour tout rentier naturel $n,\: u_n \leqslant u_{n+1} \leqslant \dfrac12$
\end{itemize}

\item 
Cet encadrement montre que 
\begin{itemize}
\item la suite $\left(u_n\right)$ est croissante et
\item la suite $\left(u_n\right)$ est majorée par $\dfrac12$.
\end{itemize}

Conclusion la suite $\left(u_n\right)$ croissante et majorée converge vers une limite $\ell \leqslant \dfrac12$.

\item 
La fonction $f$ étant continue on a donc :

$\displaystyle\lim_{n \to + \infty} f\left(u_n\right) = \displaystyle\lim_{n \to + \infty} 2u_n\left(1 - u_n\right)$, soit 

$\ell = 2\ell(1 - \ell) \iff \ell = 2\ell - 2\ell^2 \iff 2\ell^2 - \ell = 0 \iff \ell\left(2\ell - 1\right) = 0 \iff \left\{\begin{array}{l c l}
\ell&=&0\\
2\ell - 1&=&0
\end{array}\right. \iff \left\{\begin{array}{l c l}
\ell&=&0\\
\ell &=&\dfrac12
\end{array}\right.$ La seule solution acceptable est $\ell = \dfrac12$.
\end{enumerate}

\bigskip

\textbf{Partie B}

\medskip











\begin{center}$P_{n+1} -  P_n = P_n\left(1 -  b \times P_n\right)$, où $b$ est un réel strictement positif.\end{center}

Le réel $b$ est un facteur de freinage qui permet de tenir compte du caractère limité des
ressources du milieu dans lequel évoluent ces individus.

\medskip

\begin{enumerate}
\item Dans cette question $b  = 0$.
	\begin{enumerate}
		\item 
		La relation de récurrence s'écrit alors $P_{n+1} - P_n = P_n\left(1 - 0 \times P_n\right)$, soit 

$P_{n+1} - P_n = P_n \iff P_{n+1}  = 2 P_n$, donc $P_{n+1} = 2 P_n$ : cette égalité montre que la suite $\left(P_n\right)$ est une suite géométrique de raison $2$ et de premier terme $P_0 = 3$.
		\item 
On sait quel que soit $n \in \N$, \: $P_n = 3 \times 2^n$ : la limite de la suite $\left(P_n\right)$ est plus l'infini (impossible)
	\end{enumerate}
\item Dans cette question $b = 0,2$. On a donc $P_{n+1} - P_n = P_n\left(1 - 0,2P_n\right)$
	\begin{enumerate}
		\item 


On a $v_0 = 0,1 \times P_0 = 01 \times 3 = 0,3$.

La relation de récurrence devient :

$P_{n+1} -  P_n = P_n\left(1 -  0,2 \times P_n\right) \iff P_{n+1} -  P_n = P_n - 0,2P_n^2$, d'où

$P_{n+1} = 2P_n - 0,2P_n^2$. \: (1)

Or $v_{n+1} = 0,1P_{n+1} = 0,1\left(2P_n - 0,2P_n^2\right) = 0,2\left(P_n - 0,1P_n^2\right)$.

Comme $v_n = 0,1 \times P_n \iff 10v_n = P_n$, l'égalité ci-dessus devient :

$v_{n+1} = 0,2\left(10v_n - 0,1 \times \left(10v_n\right)^2\right)$, soit 

$v_{n+1} = 0,2\left(10v_n - 10v_n^2\right) = 2v_n \left(1 - v_n\right)$.
		\item 
La relation de la question précédente montre que la suite $\left(v_n \right)$ est la suite $\left(u_n \right)$ de la \textbf{Partie A} et on a vu que cette suite converge vers $\dfrac12$.

On a donc $\displaystyle\lim_{n \to + \infty} v_n = \dfrac12$.

Or $P_n = 10v_n$, donc $\displaystyle\lim_{n \to + \infty}P_n = 10 \times \dfrac12 = 5$.

Finalement la population se rapprochera de \np{5000}~individus.
	\end{enumerate}
\end{enumerate}

\end{document}
